El precedente clásico es la adaptación de discretizaciones mediante estimadores de error. FEM permite incorporar resolución local de forma explícita \cite{dorfler1996}. En aprendizaje informado por física, la incorporación de ecuaciones a una pérdida permitió tratar problemas directos e inversos con aproximadores neuronales \cite{lagaris1998,raissi2019}.

El desarrollo posterior incluye muestreo residual, formulaciones variacionales y representación por subdominios. KAN y PIKAN añaden funciones aprendibles sobre aristas; la extensión de grilla forma parte de esa literatura desde su origen \cite{liu2024kan}. La investigación local en vigas piezoeléctricas aporta un caso acoplado y una referencia de formulación \cite{gonzalez2026pinnpiezo}.

La comparación técnica de estrategias se concentra en la sección 6.3. Este capítulo establece la genealogía del problema sin duplicar ecuaciones ni el inventario del estado del arte. Los preprints recientes se incorporan como trabajos públicos que informan la dirección y las pruebas, sin suponer que sus resultados hayan sido reproducidos en este proyecto.
